%==============================================================================
\section{Required Skills --- Các kỹ năng cần có cho học máy}
%==============================================================================

\begin{dongco}[title={Tổng quan}]
Học máy là lĩnh vực phát triển nhanh và để thành công thường cần kết hợp \textbf{kỹ năng kỹ thuật} và \textbf{kỹ năng mềm}.
Khuyến nghị: nếu muốn theo học máy như một hướng nghề nghiệp, hãy học đủ các kỹ năng giúp tăng năng lực tổng thể.
\end{dongco}

\subsection{Danh sách kỹ năng chính}

\begin{phuongphap}[title={Key skills}]
\begin{itemize}
  \item Kỹ năng lập trình (Programming Skills)
  \item Thống kê và Toán (Statistics and Mathematics)
  \item Cấu trúc dữ liệu (Data Structures)
  \item Tiền xử lý dữ liệu (Data Preprocessing)
  \item Trực quan hoá dữ liệu (Data Visualization)
  \item Thuật toán học máy (Machine Learning Algorithms)
  \item Mạng nơ-ron và Học sâu (Neural Networks \& Deep Learning)
  \item Xử lý ngôn ngữ tự nhiên (Natural Language Processing)
  \item Kỹ năng giải quyết vấn đề (Problem-solving Skills)
  \item Kỹ năng giao tiếp (Communication Skills)
  \item Hiểu kinh doanh (Business Acumen)
\end{itemize}
\end{phuongphap}

\subsection{Programming Skills}

\begin{ynghia}[title={Vì sao cần lập trình?}]
Học máy cần nền tảng lập trình vững, đặc biệt ở các ngôn ngữ như Python, R và Java.
Thành thạo lập trình giúp bạn xây, kiểm thử và triển khai mô hình ML.
\end{ynghia}

\begin{phuongphap}[title={Python và R}]
\begin{itemize}
  \item Python rất phổ biến vì có nhiều thư viện như NumPy, Matplotlib, scikit-learn, Seaborn, Keras, TensorFlow, \ldots giúp đơn giản hoá quy trình.
  \item Các nội dung Python cơ bản: basic data types; dictionaries/lists/sets; loops/conditionals; functions; list comprehensions.
  \item R cũng phổ biến trong ML; có thể làm các tác vụ ML nặng “dễ hơn”; ngoài ngôn ngữ, cần nắm các package đi kèm.
\end{itemize}
\end{phuongphap}

\subsection{Statistics and Mathematics}

\begin{dongco}[title={Toán--thống kê là nền để hiểu mô hình}]
Hiểu tốt toán và thống kê giúp bạn hiểu/áp dụng mô hình thống kê, thuật toán và phương pháp để phân tích và diễn giải dữ liệu.
\end{dongco}

\subsubsection{Statistics}
\begin{ynghia}[title={Thống kê dùng để làm gì?}]
Thống kê giúp suy luận từ dữ liệu và rút kết luận.
Công thức thống kê giúp diễn giải dữ liệu để ra quyết định dựa trên dữ liệu.
Chia thống kê thành:
\begin{itemize}
  \item \textbf{Thống kê mô tả (Descriptive statistics)}: đơn giản hoá/tổ chức dữ liệu bằng mean, range, variance, standard deviation.
  \item \textbf{Thống kê suy luận (Inferential statistics)}: dùng mẫu nhỏ để kết luận về tập lớn, qua hypothesis testing, null/alternative testing, \ldots
\end{itemize}
\end{ynghia}

\subsubsection{Mathematics}
\begin{phuongphap}[title={Các mảng toán nên biết}]
\begin{itemize}
  \item \textbf{Đại số (Algebra)}: biến, hằng, hàm; phương trình tuyến tính; logarit.
  \item \textbf{Đại số tuyến tính (Linear algebra)}: vectơ, ma trận, trị riêng (eigenvalues).
  \item \textbf{Giải tích (Calculus)}: đạo hàm, tích phân, gradient descent.
\end{itemize}
\end{phuongphap}

\begin{vidu}[title={Toán liên quan ML như thế nào?}]
Ví dụ: công thức hồi quy tuyến tính \(y=ax+b\) là biểu thức tuyến tính trong đại số.
\end{vidu}

\subsubsection{Mathematical Notation}
\begin{luuy}[title={Đọc được ký hiệu quan trọng hơn “làm tay phức tạp”}]
Mức toán cần biết có thể chỉ ở mức beginner, nhưng điều quan trọng là bạn đọc và hiểu ký hiệu toán trong phương trình.
Nếu đọc và hiểu được ký hiệu, bạn sẵn sàng học ML; nếu không, nên ôn lại toán.
\end{luuy}

\subsection{Probability Theory}

\begin{ynghia}[title={Xác suất là tiền đề quan trọng}]
Vì ML liên quan dự đoán, xác suất là nền tảng quan trọng.
Gợi ý nắm: biến ngẫu nhiên, mật độ/phân phối xác suất, \ldots
\end{ynghia}

\begin{vidu}[title={Ví dụ kiểm tra kiến thức: phân loại theo xác suất có điều kiện}]
\[
p(c_{i}|x,y)=\frac{p(x,y|c_{i})\,p(c_{i})}{p(x,y)}.
\]
Quy tắc phân loại Bayes:
\begin{itemize}
  \item Nếu \(P(c_1|x,y) > P(c_2|x,y)\) thì lớp là \(c_1\).
  \item Nếu \(P(c_1|x,y) < P(c_2|x,y)\) thì lớp là \(c_2\).
\end{itemize}
\end{vidu}

\subsection{Optimization Problem}

\begin{ynghia}[title={Tối ưu hoá trong ML}]
Ví dụ một bài toán tối ưu (mang tính kiểm tra khả năng đọc ký hiệu).
Trong ML, tối ưu hoá thường xuất hiện khi huấn luyện mô hình (tối thiểu hoá loss / tối đa hoá mục tiêu).
\end{ynghia}

\subsection{Data Structures}

\begin{ynghia}[title={Vì sao cần cấu trúc dữ liệu?}]
Hiểu cấu trúc dữ liệu giúp giải bài toán thực và xây sản phẩm phần mềm.
Trong ML, cấu trúc dữ liệu giúp xử lý và hiểu bài toán phức tạp.
Một số cấu trúc: arrays, stack, queue, binary trees, maps, \ldots
\end{ynghia}

\subsection{Data Preprocessing}

\begin{ynghia}[title={Tiền xử lý dữ liệu}]
Chuẩn bị dữ liệu cho ML cần kiến thức về làm sạch, biến đổi và chuẩn hoá dữ liệu.
Điều này gồm: phát hiện và sửa lỗi, missing values, và bất nhất trong dữ liệu.
\end{ynghia}

\subsection{Data Visualization}

\begin{ynghia}[title={Trực quan hoá dữ liệu}]
Trực quan hoá là tạo biểu diễn đồ hoạ để hiểu và diễn giải dataset phức tạp.
Data scientist cần tạo visualizations hiệu quả để truyền đạt insight.
Công cụ: Tableau, Power BI, \ldots
\end{ynghia}

\begin{luuy}[title={Hiểu các loại biểu đồ}]
Nhiều trường hợp bạn cần hiểu các loại biểu đồ để hiểu phân phối dữ liệu và diễn giải output của thuật toán.
\end{luuy}

\subsection{Machine Learning Algorithms}

\begin{ynghia}[title={Kiến thức thuật toán}]
Liệt kê: regression, decision trees, random forests, k-NN, SVM, neural networks.
Hiểu điểm mạnh/yếu của thuật toán quan trọng để xây mô hình hiệu quả và biết khi nào áp dụng.
\end{ynghia}

\subsection{Neural Networks \& Deep Learning}

\begin{ynghia}[title={NN và DL}]
Neural networks được thiết kế để máy hoạt động “tương tự não người”, gồm các nút/neurons liên kết để học từ dữ liệu.
Deep learning là nhánh của ML huấn luyện mạng sâu để phân tích dataset phức tạp; cần hiểu CNN, RNN và các chủ đề liên quan.
\end{ynghia}

\subsection{Natural Language Processing}

\begin{ynghia}[title={NLP}]
NLP là nhánh AI tập trung vào tương tác giữa máy và người bằng ngôn ngữ tự nhiên.
Kỹ thuật như sentiment analysis, text classification, named entity recognition.
\end{ynghia}

\subsection{Problem-solving Skills}

\begin{ynghia}[title={Kỹ năng giải quyết vấn đề}]
Cần khả năng xác định vấn đề, tạo giả thuyết và phát triển lời giải; cần tư duy sáng tạo và logic cho bài toán phức tạp.
\end{ynghia}

\subsection{Communication Skills}

\begin{ynghia}[title={Kỹ năng giao tiếp}]
Data scientist cần giải thích khái niệm kỹ thuật cho stakeholder không kỹ thuật, truyền đạt kết quả phân tích và ý nghĩa một cách rõ ràng, súc tích.
\end{ynghia}

\subsection{Business Acumen}

\begin{ynghia}[title={Hiểu bối cảnh kinh doanh}]
Vì ML thường dùng để giải bài toán kinh doanh, hiểu bối cảnh và biết cách áp dụng ML vào vấn đề kinh doanh là quan trọng.
\end{ynghia}

\begin{luuy}[title={Kết luận}]
Học máy cần nhiều kỹ năng (kỹ thuật, toán, kỹ năng mềm) và cần kết hợp chúng để xây dựng và triển khai mô hình hiệu quả.
\end{luuy}

\begin{vidu}[title={Hình minh hoạ kỹ năng nền tảng}]
\tsimg{04_yeu_cau_nen_tang_va_ky_nang/img01.jpg}
\tsimg{04_yeu_cau_nen_tang_va_ky_nang/img02.jpg}
\end{vidu}
